\documentclass[a4paper,12pt]{article}
\usepackage[T2A]{fontenc}
\usepackage[utf8]{inputenc}
\usepackage[russian]{babel}
\usepackage{amsmath,amssymb,mathtools}
\renewcommand{\familydefault}{\sfdefault}
\usepackage{icomma}
\usepackage[a4paper,left=23mm,right=23mm,top=24mm,bottom=24mm,headheight=15pt]{geometry}
\usepackage{microtype}
\usepackage{tikz}
\usetikzlibrary{arrows.meta,positioning,shapes.geometric,calc}
\usepackage{booktabs,tabularx}
\usepackage{enumitem}
\usepackage{fancyhdr}
\usepackage{xcolor}
\usepackage{setspace}
\usepackage{needspace}
\usepackage{hyperref}
\definecolor{accent}{HTML}{247C80}
\definecolor{darkaccent}{HTML}{13575C}
\definecolor{textgray}{HTML}{334047}
\definecolor{softgray}{HTML}{617178}
\definecolor{paleblue}{HTML}{E8F2F1}
\hypersetup{colorlinks=true,linkcolor=darkaccent}
\setlength{\parindent}{0pt}
\setlength{\parskip}{4pt}
\onehalfspacing
\setlist[enumerate]{leftmargin=*,itemsep=10pt,topsep=8pt}
\pagestyle{fancy}
\fancyhf{}
\fancyhead[L]{\small\color{softgray} ЕГЭ · математика}
\fancyhead[R]{\small\color{softgray} 16.10.2026}
\fancyfoot[L]{\footnotesize\color{softgray} Николь Саркисьянц}
\fancyfoot[R]{\footnotesize\color{softgray}\thepage}
\renewcommand{\headrulewidth}{0.25pt}
\renewcommand{\footrulewidth}{0pt}
\newcommand{\sectionnote}[1]{{\small\color{softgray}#1\par}}
\newcommand{\criterion}[1]{\textcolor{darkaccent}{\bfseries #1}}
\newcommand{\checkitem}[1]{\par\noindent\hspace{-1pt}\textcolor{accent}{\(\square\)}\enspace #1\par}
\newcommand{\lessonsection}[1]{\Needspace{5\baselineskip}\section*{\color{darkaccent}#1}\vspace{-5pt}\noindent\textcolor{accent}{\rule{25mm}{0.7pt}}\par\vspace{1pt}}
\begin{document}
\thispagestyle{empty}
\begin{titlepage}
\begin{center}
\vspace*{29mm}
{\large\color{darkaccent}\bfseries МАТЕМАТИКА · ЕГЭ}\par
\vspace{28mm}
{\Huge\bfseries Чек-лист}\par
\vspace{9mm}
{\LARGE\color{darkaccent}\bfseries Финансы и тригонометрия}\par
\vspace{4mm}
{\Large Повторение и работа над ошибками}\par
\vspace{21mm}
\textcolor{accent}{\rule{34mm}{0.7pt}}\par
\vspace{15mm}
{\large Николь Саркисьянц}\par
\vspace{6mm}
{\large 16 октября 2026 года}\par
\vfill
{\small\color{textgray}Лёвин Артём Александрович}\par
{\small\color{textgray}эксклюзивно для Николь Саркисьянц}\par
\vspace{17mm}
\end{center}
\begin{tikzpicture}[remember picture,overlay]
\draw[accent!55,line width=1.1pt] ($(current page.south west)+(0,16mm)$) .. controls ($(current page.south west)+(53mm,29mm)$) and ($(current page.south east)+(-72mm,3mm)$) .. ($(current page.south east)+(0,14mm)$);
\end{tikzpicture}
\end{titlepage}

\lessonsection{01. Кредит с равномерным уменьшением долга}
\sectionnote{Дифференцированная схема: основной долг погашается равными долями.}

\checkitem{\criterion{Узнаю схему по условию:} долг в одну и ту же дату каждого года (или месяца) становится на одну и ту же величину меньше.}
\checkitem{\criterion{Разделяю понятия:} основной долг, проценты за период, платёж и остаток долга.}
\checkitem{\criterion{Перевожу проценты в коэффициент:} \(p=\dfrac r{100}\), где \(r\)~--- ставка в процентах за период.}

Пусть \(S\)~--- сумма кредита, \(n\)~--- число одинаковых периодов, \(p\)~--- процентная ставка за один период в долях единицы. Проценты начисляются \emph{до платежа} на остаток долга предыдущего периода.

\begin{align*}
\underbrace{\frac Sn}_{\text{доля основного долга}}&\quad\text{погашается каждый период},\\
D_i&=S\left(1-\frac{i-1}{n}\right) &&\text{долг перед начислением процентов},\\
P_i&=\frac Sn+pD_i &&\text{платёж в период }i,\quad i=1,\ldots,n.
\end{align*}

\checkitem{\criterion{Знаю основные формулы} и условия их применимости.}
\begin{align*}
\text{Общая сумма платежей}\quad A&=S+\frac{Sp(n+1)}{2},\\
\text{Переплата}\quad \Delta&=\frac{Sp(n+1)}{2},\\
\text{Переплата в процентах}\quad q&=\frac{p(n+1)}{2}\cdot100\%,\\
\text{Первый (максимальный) платёж}\quad P_1&=\frac Sn+Sp,\\
\text{Последний (минимальный) платёж}\quad P_n&=\frac Sn+\frac{Sp}{n}.
\end{align*}
\sectionnote{Для положительной ставки \(p>0\) платежи строго уменьшаются. При \(p=0\) они равны.}

\clearpage
\lessonsection{02. Как из условия получить ставку}
\checkitem{\criterion{Умею составить таблицу} для двух периодов, чтобы видеть происхождение формул.}

\textbf{Пример.} Кредит \(1400\) тыс. руб. на два года. Общая сумма выплат \(1736\) тыс. руб. Найдите годовую процентную ставку.

\renewcommand{\arraystretch}{1.24}
\begin{center}
\small
\begin{tabularx}{\linewidth}{@{}cXXXX@{}}
\toprule
Год & Долг перед начислением & Проценты & Платёж & Остаток \\
\midrule
1 & \(1400\) & \(1400p\) & \(700+1400p\) & \(700\) \\
2 & \(700\) & \(700p\) & \(700+700p\) & \(0\) \\
\bottomrule
\end{tabularx}
\end{center}
\[
(700+1400p)+(700+700p)=1736,\qquad
2100p=336.
\]
\[
p=0{,}16,\qquad \boxed{r=16\%}.
\]
\checkitem{\criterion{Умею обойтись без таблицы:} для тех же данных \(1400+\dfrac{1400p\cdot3}{2}=1736\).}

\textbf{Ещё один пример.} За \(14\) месяцев общая сумма выплат оказалась на \(15\%\) больше суммы кредита. При ежемесячном уменьшении основного долга на равную величину
\[
\frac{p(14+1)}2=0{,}15,\quad p=0{,}02,\quad \boxed{r=2\%\text{ в месяц}}.
\]

\checkitem{\criterion{Различаю} \(15\%\) переплаты за весь срок и \(15\%\) процентной ставки за один период.}
\checkitem{\criterion{Помню размерность:} годы и годовая ставка; месяцы и месячная ставка.}

\lessonsection{03. Ограничения на платежи}
\checkitem{\criterion{Перевожу слова в знаки:} «не более»~--- \(\le\), «не менее»~--- \(\ge\).}
\checkitem{\criterion{Использую систему неравенств} для наибольшего и наименьшего платежей.}

\textbf{Пример.} Кредит \(6\) млн руб. на \(15\) лет. Наибольший платёж \emph{не более} \(1{,}9\) млн, наименьший~--- \emph{не менее} \(0{,}5\) млн. Найдите ставку.
\[
\begin{cases}
\dfrac6{15}+6p\le 1{,}9,\\[5pt]
\dfrac6{15}+\dfrac6{15}p\ge0{,}5
\end{cases}
\iff
\begin{cases}
p\le0{,}25,\\p\ge0{,}25.
\end{cases}
\]
Следовательно, \(p=0{,}25\), то есть \(\boxed{r=25\%}\). Оба граничных платежа здесь равны заданным порогам.

\clearpage
\thispagestyle{fancy}
\begin{center}
{\Large\color{darkaccent}\bfseries Алгоритм: дифференцированный кредит}\par
\vspace{3mm}
\textcolor{softgray}{От условия задачи к нужной формуле}\par
\vspace{9mm}
\end{center}
\tikzset{
 algterm/.style={ellipse,draw=darkaccent,very thick,align=center,text width=6.2cm,minimum height=11mm,font=\small},
 algstep/.style={rectangle,rounded corners=2mm,draw=accent,line width=0.8pt,align=center,text width=9.2cm,minimum height=15mm,font=\small,inner sep=4pt},
 algtest/.style={diamond,aspect=2.7,draw=darkaccent,line width=0.85pt,align=center,text width=5.6cm,inner sep=2pt,font=\small},
 algarrow/.style={-{Stealth[length=3mm,width=2mm]},line width=0.85pt,draw=darkaccent}
}
\begin{center}
\begin{tikzpicture}[node distance=9mm]
\node[algterm] (start) {Прочитайте условие};
\node[algtest,below=of start] (test) {Основной долг уменьшается\\на равную величину?};
\node[algstep,below=14mm of test] (data) {Определите \(S\), \(n\), \(p=r/100\).\\Приведите ставку и срок к одному периоду};
\node[algstep,below=of data] (principal) {Основная часть каждого платежа: \(S/n\).\\Долг перед периодом \(i\): \(D_i=S(1-(i-1)/n)\)};
\node[algstep,below=of principal] (formula) {Выберите формулу по вопросу: \\ \(P_i=S/n+pD_i\),\quad \(A=S+Sp(n+1)/2\)};
\node[algstep,below=of formula] (check) {Подставьте данные, решите уравнение\\или систему неравенств; проверьте единицы};
\node[algterm,below=of check] (done) {Ответ: ставка в процентах,\\сумма в нужных единицах};
\node[algstep,right=16mm of test,text width=3.1cm,minimum height=15mm] (other) {Используйте другую схему погашения};
\draw[algarrow] (start)--(test);
\draw[algarrow] (test)--node[right,font=\footnotesize]{да}(data);
\draw[algarrow] (test.east)--node[above,font=\footnotesize]{нет}(other.west);
\draw[algarrow] (data)--(principal);
\draw[algarrow] (principal)--(formula);
\draw[algarrow] (formula)--(check);
\draw[algarrow] (check)--(done);
\end{tikzpicture}
\end{center}
\clearpage

\lessonsection{04. Тригонометрия: внимание к преобразованиям}
\checkitem{\criterion{Делю все части равенства корректно.} При делении на \(\cos x\) каждый член делится на \(\cos x\), а случай \(\cos x=0\) нужно рассмотреть отдельно.}
\[
\sin x+\cos x=1
\quad\Longrightarrow\quad
\tg x+1=\frac1{\cos x}\quad (\cos x\ne0).
\]
\sectionnote{Заменять правую часть на \(1\) после деления нельзя. Деление на \(\cos x\) особенно удобно в однородном уравнении с правой частью \(0\), после проверки нулевых случаев.}

\checkitem{\criterion{Помню признаки арифметической ошибки:} \(-3+\frac12=\frac{-6+1}{2}=-\frac52\), а \(24+20=44\).}
\checkitem{\criterion{Контролирую вынесение минуса:} \(\dfrac{-6+1}{2}=-\dfrac{6-1}{2}\).}
\checkitem{\criterion{Проверяю, что после деления на выражение не потеряны корни, при которых оно обращается в нуль.}}

\lessonsection{05. Формулы понижения степени}
\checkitem{\criterion{Знаю формулы} для синуса и косинуса.}
\[
\boxed{\sin^2u=\frac{1-\cos2u}{2}},\qquad
\boxed{\cos^2u=\frac{1+\cos2u}{2}}.
\]
\checkitem{\criterion{Понимаю ограничения} для тангенса и котангенса.}
\[
\tg^2u=\frac{1-\cos2u}{1+\cos2u}\quad(\cos u\ne0),
\]
\[
\ctg^2u=\frac{1+\cos2u}{1-\cos2u}\quad(\sin u\ne0).
\]
\checkitem{\criterion{Применяю формулу сразу к аргументу:}}
\[\cos^2(2x)=\frac{1+\cos4x}{2},\qquad \sin^2(5x)=\frac{1-\cos10x}{2}.\]
\checkitem{\criterion{Вспоминаю простейшее уравнение:} \(\cos u=0\iff u=\dfrac\pi2+\pi k,\ k\in\mathbb Z\).}

\clearpage
\lessonsection{06. Показательно-тригонометрическое уравнение}
\checkitem{\criterion{Замечаю повторяющееся выражение} и делаю замену, сохраняя диапазон значений.}
\checkitem{\criterion{Привожу степени к общему показателю} и выношу общий множитель.}

\textbf{Пример.} Решите уравнение
\[
4^{2+\cos^2 2x}-4^{\cos^2 2x}
=9^{1+\cos^2 2x}+9^{\cos^2 2x}.
\]
\textbf{Шаг 1.} Положим \(t=\cos^2 2x\), где \(0\le t\le1\):
\[
4^{2+t}-4^t=9^{1+t}+9^t.
\]
\textbf{Шаг 2.} Вынесем общие степени:
\[
4^t(16-1)=9^t(9+1)
\quad\Longleftrightarrow\quad 15\cdot4^t=10\cdot9^t.
\]
\textbf{Шаг 3.} Сведём к одному основанию:
\[
\left(\frac49\right)^t=\frac23
\quad\Longleftrightarrow\quad
\left(\frac23\right)^{2t}=\left(\frac23\right)^1
\quad\Longleftrightarrow\quad t=\frac12.
\]
\textbf{Шаг 4.} Вернём замену и понизим степень:
\[
\cos^2 2x=\frac12
\iff \frac{1+\cos4x}{2}=\frac12
\iff \cos4x=0.
\]
\[
4x=\frac\pi2+\pi k
\quad\Longrightarrow\quad
\boxed{x=\frac\pi8+\frac{\pi k}{4}},\qquad k\in\mathbb Z.
\]
\checkitem{\criterion{Проверяю обратную подстановку:} \(t=\frac12\) допустимо, так как \(0\le t\le1\).}
\checkitem{\criterion{Выбираю более короткое решение:} понижение степени здесь позволяет обойтись одним тригонометрическим уравнением.}

\clearpage
\lessonsection{07. Самостоятельная тренировка · Путь А}
\sectionnote{10 заданий, от базовых к более содержательным. Решения оформляйте самостоятельно; ответы помещены после тренировочного блока.}
\begin{enumerate}[label=\textbf{\arabic*.},itemsep=12pt]
\item Кредит \(1200\) тыс. руб. взят на два года под \(10\%\) годовых; основной долг уменьшается равномерно. Найдите оба годовых платежа и их сумму.
\item Кредит \(1500\) тыс. руб. взят на четыре года под \(8\%\) годовых по дифференцированной схеме. Найдите общую сумму выплат и переплату.
\item По кредиту \(2000\) тыс. руб. на пять лет с равномерным погашением основного долга выплатили всего \(2300\) тыс. руб. Найдите годовую процентную ставку.
\item Кредит с равномерным погашением основного долга взят на \(14\) месяцев. Общая сумма платежей превышает сумму кредита на \(15\%\). Найдите месячную процентную ставку.
\item Кредит \(8\) млн руб. взят на \(20\) лет. Каждый год основной долг уменьшается на одинаковую величину. Первый годовой платёж \emph{не более} \(1{,}2\) млн руб., последний~--- \emph{не менее} \(0{,}44\) млн руб. Найдите годовую процентную ставку.
\item Решите уравнение \(\tg\dfrac x2-1=0\).
\item Решите уравнение \(\sin^2(5x)=\dfrac12\), используя понижение степени.
\item Решите уравнение \(\cos^2(3x)=\dfrac34\), используя понижение степени.
\item Решите уравнение \(4^{2+\cos^2 x}-4^{\cos^2 x}=9^{1+\cos^2 x}+9^{\cos^2 x}\).
\item Решите уравнение \(\sin x+\cos x=1\). Объясните, почему прямое деление обеих частей на \(\cos x\) требует отдельной проверки случая \(\cos x=0\).
\end{enumerate}

\clearpage
\lessonsection{08. Ответы для самопроверки}
\sectionnote{В финансовых заданиях ставки указаны за один период, соответствующий условию. \(k\in\mathbb Z\).}
\renewcommand{\arraystretch}{1.65}
\begin{tabularx}{\linewidth}{@{}p{13mm}X@{}}
\toprule
\textbf{№} & \textbf{Ответ}\\
\midrule
1 & \(720\) тыс. руб.; \(660\) тыс. руб.; всего \(1380\) тыс. руб.\\
2 & \(1800\) тыс. руб.; переплата \(300\) тыс. руб.\\
3 & \(5\%\) годовых.\\
4 & \(2\%\) в месяц.\\
5 & \(10\%\) годовых.\\
6 & \(x=\dfrac\pi2+2\pi k\).\\
7 & \(x=\dfrac\pi{20}+\dfrac{\pi k}{10}\).\\
8 & \(x=\pm\dfrac\pi{18}+\dfrac{\pi k}{3}\).\\
9 & \(x=\dfrac\pi4+\dfrac{\pi k}{2}\).\\
10 & \(x=2\pi k\) или \(x=\dfrac\pi2+2\pi k\).\\
\bottomrule
\end{tabularx}

\vspace{11mm}
\textbf{Контроль перед сдачей}
\checkitem{В финансовой задаче я проверил(-а), соответствует ли ставка периоду.}
\checkitem{В тригонометрическом уравнении я проверил(-а) все деления и ограничения.}
\checkitem{В ответе я указал(-а) все серии корней и параметр \(k\in\mathbb Z\).}
\end{document}
